\documentclass[11pt]{article}

\usepackage[a4paper,margin=1in]{geometry}
\usepackage{amsmath}
\usepackage{amssymb}
\usepackage{booktabs}
\usepackage{hyperref}
\usepackage{xcolor}
\usepackage{listings}

\lstset{
  basicstyle=\ttfamily\small,
  breaklines=true,
  frame=single,
  columns=fullflexible
}

\title{Robot + Object MCGDF World Model:\\Structured Contact via Newton's Third Law and Newton--Euler Object Dynamics}
\author{}
\date{}

\begin{document}
\maketitle

\section{Overview}

The implementation under
\begin{center}
\texttt{scripts/world\_model/Physics/Robot\_and\_Object/MCGDF}
\end{center}
extends the robot-only MCGDF model with a structured robot--object coupling.
It addresses the mismatches identified in Section~4.3 of the verification
report
\texttt{scripts/world\_model/Physics/delan\_vs\_physx\_verification.tex},
in particular the structural items \#1 (no action--reaction coupling), \#5
(no gyroscopic object structure), \#6 (body-frame inertia in world-frame
dynamics), \#7 (linearly dependent object linear-acceleration targets),
\#9 (env-frame object position), and Approach~A of \#2 (decompose the
residual target).

\section{State, Input, and Context}

The state is
\begin{equation}
  s_t =
  \begin{bmatrix}
    q_t \\ \dot q_t \\ p^o_t \\ r^o_t \\ v^o_t \\ \omega^o_t
  \end{bmatrix} \in \mathbb{R}^{18 + 13} = \mathbb{R}^{31},
\end{equation}
the input is the recorded torque $\tau_t \in \mathbb{R}^9$, and the
per-episode physical context is
\begin{equation}
  \xi = \bigl[\,m_o,\ I_b\,\mathrm{(9 entries)},\ \mu_o\,(3)\,\bigr] \in \mathbb{R}^{13}.
\end{equation}
$I_b$ is the cube inertia in its body frame.

\section{Forward Dynamics}

\subsection{Robot side}
\begin{equation}
  H(q)\,\ddot q
  = \tau + r_\theta(q,\dot q,\tau)
  - c(q,\dot q) - g(q)
  - d\dot q - f\tanh(\dot q/\varepsilon_f)
  - J_c(q)^\top F_c(s,\tau,\xi),
  \label{eq:mcgdf-robot-coupled}
\end{equation}
where $J_c \in \mathbb{R}^{6\times n}$ is the analytic geometric Jacobian of
the Franka gripper, evaluated at $q$ (last $n-7$ columns are zero), and
$F_c \in \mathbb{R}^6$ is the learned contact wrench (linear + angular) on
the cube in world frame. The robot-side contact torque
$-J_c(q)^\top F_c$ enforces Newton's third law structurally.

Damping is omitted from the RHS by default (\texttt{--omit\_damping}) for
the same reason as in the robot-only MCGDF: under Franka's implicit-actuator
configuration, the recorded \texttt{applied\_torque} already contains the
implicit-PD damping torque.

\subsection{Object side}
\begin{align}
  m_o\,\dot v^o
  &= m_o\, g_w + F_c^{\mathrm{lin}},
  \label{eq:mcgdf-object-linear}\\
  I_w(r^o)\,\dot\omega^o + \omega^o\times\bigl(I_w(r^o)\,\omega^o\bigr)
  &= F_c^{\mathrm{ang}},
  \label{eq:mcgdf-object-angular}\\
  I_w(r^o) &= R(r^o)\,I_b\,R(r^o)^\top.
  \label{eq:mcgdf-Iw}
\end{align}
The world-frame inertia $I_w$ is constructed every step from the
predicted/recorded quaternion and the body-frame inertia tensor stored in
$\xi$, fixing mismatch~\#6. The gyroscopic term $\omega^o\times(I_w\omega^o)$
is part of the equation rather than learned (mismatch~\#5).

\subsection{Quaternion integration}
The object orientation is integrated by the closed-form exponential map
\begin{equation}
  r^o_{t+1} = r^o_t \otimes
  \exp\!\bigl(\tfrac{1}{2}\,\omega^o_{t+1}\,\Delta t\bigr),
  \quad
  \exp\!\bigl([0,\mathbf{v}]\bigr) = \bigl(\cos\|\mathbf{v}\|,\,\tfrac{\sin\|\mathbf{v}\|}{\|\mathbf{v}\|}\mathbf{v}\bigr),
  \label{eq:mcgdf-quat-exp}
\end{equation}
replacing the linear-Euler step used in the GT-dynamics variant
(mismatch~\#4).  Position and velocities use semi-implicit Euler.

\section{Neural Network Architecture}
\label{sec:architecture}

The model contains five focused sub-networks plus optional 18 scalar
parameters when joint damping/friction are made learnable. Everything else
in the forward pass (forward kinematics, Newton--Euler object dynamics,
quaternion integration, semi-implicit Euler) is analytic.

\subsection{Sub-network inventory}

\begin{table}[h]
\centering
\small
\begin{tabular}{p{0.27\linewidth} p{0.22\linewidth} p{0.18\linewidth} p{0.11\linewidth} p{0.16\linewidth}}
\toprule
Sub-network (class)              & Input                                  & Output                       & Depth & Hidden \\
\midrule
\texttt{g\_net} (MLP)            & $q \in \mathbb{R}^9$                  & $g \in \mathbb{R}^9$         & 2     & 256    \\
\texttt{l\_diag\_net} (Deriv.\,MLP) & $q \in \mathbb{R}^9$              & $\ell_d \in \mathbb{R}^9$    & 2     & 256    \\
\texttt{l\_offdiag\_net} (Deriv.\,MLP) & $q \in \mathbb{R}^9$           & $\ell_o \in \mathbb{R}^{36}$ & 2     & 256    \\
\texttt{ResidualTorqueHead}      & $(q,\dot q,\tau) \in \mathbb{R}^{27}$ & $r_\theta \in \mathbb{R}^9$  & 2     & 128    \\
\texttt{ContactWrenchHead}       & $(s,\tau,\xi) \in \mathbb{R}^{53}$    & $F_c \in \mathbb{R}^6$       & 3     & 128    \\
\texttt{JointDampingFrictionParams} & ---                                & 9+9 scalars (softplus, opt.) & --- & ---    \\
\bottomrule
\end{tabular}
\caption{Trainable sub-networks of the robot+object MCGDF model.
``Depth'' is the number of hidden layers, so each row has $d+1$ linear
layers separated by SiLU activations. All sizes are exposed as CLI flags.}
\label{tab:subnets}
\end{table}

\subsection{What each sub-network produces}

\paragraph{(1) DeLaN structured terms.}
The robot inertia matrix $H(q)$ is not a direct MLP output. The two
inertia heads produce a Cholesky-style lower-triangular factor $L(q)$:
\texttt{l\_diag\_net} outputs 9 raw values passed through
$\mathrm{softplus}(\cdot)+\varepsilon_d$ to guarantee a positive diagonal,
\texttt{l\_offdiag\_net} outputs 36 raw values that fill the strict
lower-triangular entries. Both heads are \texttt{DerivativeMLP}s that
analytically propagate $\partial L/\partial q$ through the SiLU+Linear
chain without invoking autograd. The structural inertia and Coriolis terms
are then
\begin{equation}
  H(q) = L(q)L(q)^\top + \varepsilon I,
  \qquad
  c(q,\dot q) = \dot H \dot q - \tfrac{1}{2}\nabla_q\bigl(\dot q^\top H \dot q\bigr),
\end{equation}
with $c$ derived analytically from $L$, $\partial L/\partial q$, and
$\dot q$ (the DeLaN identity). The gravity vector is the direct output of
\texttt{g\_net(q)}.

\paragraph{(2) Residual torque head.}
\texttt{ResidualTorqueHead} is a small 2-layer MLP
\(r_\theta(q,\dot q,\tau) \in \mathbb{R}^9\)
that absorbs the implicit-PD bias and finite-difference / numerical residual
identified in the verification report.  Crucially the recorded torque
$\tau$ is included as a network input, because $\tau$ encodes the PD targets
that drove the implicit actuator.

\paragraph{(3) Contact wrench head.}
\texttt{ContactWrenchHead} is a 3-layer MLP
\begin{equation}
  F_c(s,\tau,\xi)
  \;=\;\bigl(F_c^{\mathrm{lin}},\ F_c^{\mathrm{ang}}\bigr) \in \mathbb{R}^{3}\!\oplus\mathbb{R}^{3}.
\end{equation}
Made deeper than $r_\theta$ because the dependence of the contact force on
gripper pose, cube state and material properties is more nonlinear than the
per-joint actuator bias.

\paragraph{(4) Joint damping/friction parameters.}
\texttt{JointDampingFrictionParams} holds at most $9 + 9$ trainable
scalars, parameterised as \texttt{softplus(raw)} so they stay non-negative.
Active only when \texttt{--learn\_damping\_friction} is passed; by default
the model reads $d$ and $f$ from the dataset.

\subsection{What is \emph{not} learned}

\begin{itemize}
  \item \textbf{Franka FK and Jacobian.} \texttt{FrankaForwardKinematics}
    encodes the seven-link kinematic chain as fixed
    \texttt{register\_buffer}s. It produces $p_{\mathrm{ee}}, R_{\mathrm{ee}},
    J_c \in \mathbb{R}^{6\times 9}$ as differentiable functions of $q$, with
    zero learnable parameters.
  \item \textbf{Newton--Euler object dynamics.}
    Eqs.~\eqref{eq:mcgdf-object-linear}--\eqref{eq:mcgdf-object-angular} are
    written out directly. The world-frame inertia
    $I_w = R(r^o)\,I_b\,R(r^o)^\top$ is rebuilt every step from the predicted
    quaternion and the body-frame inertia in $\xi$.
  \item \textbf{Quaternion exponential map.} Eq.~\eqref{eq:mcgdf-quat-exp}.
  \item \textbf{Semi-implicit Euler integration} for
    $q, \dot q, p^o, v^o, \omega^o$.
\end{itemize}

The learned content of the model is therefore restricted to the
free-space rigid-body terms ($H, c, g$), the implicit-actuator bias
$r_\theta$, and the contact wrench $F_c$. Everything coupling these
quantities into a trajectory is the equations of motion expressed in
PyTorch.

\subsection{End-to-end data flow per step}

\begin{lstlisting}
state s = [q (9), qdot (9), p_o (3), r_o (4 quat), v_o (3), omega_o (3)]   = 31
input  tau (9)
ctx    xi = [m_o (1), I_b (9 flat), mu_o (3)]                              = 13

Structural rigid-body:                              (DeLaN)
  g_net(q)         -> g (9)
  l_diag_net(q)    -> ell_d (9)
  l_offdiag_net(q) -> ell_o (36)
  -> assemble L (9x9), then H = L L^T + eps I
  -> coriolis c(q,qdot) analytic from L, dL/dq, qdot

Implicit-PD bias:                                   (Residual head)
  r_theta(q, qdot, tau) -> 9

Forward kinematics:                                 (analytic, fixed)
  FK(q) -> p_ee (3), R_ee (3x3)
  geometric Jacobian J_c (6 x 9), last 2 cols zero

Contact wrench:                                     (Contact head)
  F_c(s, tau, xi) -> F_lin (3), F_ang (3)
  tau_contact = -J_c^T F_c -> (9)

Robot side:
  H qdd = tau + r_theta - c - g - d*qdot - f*tanh(qdot/eps) + tau_contact
  qdd  = torch.linalg.solve(H, RHS)

Object side (Newton-Euler):                         (analytic)
  I_w        = R(r_o) I_b R(r_o)^T
  vdot_o     = [0,0,-9.81] + F_lin / m_o
  omegadot_o = I_w^{-1} (F_ang - omega_o x I_w omega_o)

Integration:                                        (analytic)
  semi-implicit Euler for q, qdot, p_o, v_o, omega_o
  exponential-map quaternion update for r_o
\end{lstlisting}

\subsection{Object trajectory prediction: a recursive rollout, not a trajectory decoder}
\label{sec:object-trajectory-prediction}

The previous subsection described one step of the
\texttt{RobotObjectMCGDFStep} forward dynamics.  A predicted \emph{object
trajectory} over a horizon $H$ is produced by chaining $H$ such steps in
the open-loop rollout implemented by
\texttt{MultiStepRobotObjectMCGDFWorldModel.forward}.  Given an initial
state $s_t$ (taken from the last entry of the history window), the per-
episode context $\xi$ (and latent $z$ when the context encoder is on),
and the recorded future torque sequence $\tau_{t+1:t+H}$, the model
iterates
\begin{equation}
  s_{t+h+1} \;=\; \mathrm{RobotObjectMCGDFStep}\bigl(s_{t+h},\,\tau_{t+h+1},\,\xi,\,z\bigr),
  \qquad h = 0, 1, \dots, H-1,
  \label{eq:trajectory-recursion}
\end{equation}
and the predicted object trajectory is the projection
\begin{equation}
  \bigl\{(p^o_{t+h+1},\ r^o_{t+h+1},\ v^o_{t+h+1},\ \omega^o_{t+h+1})\bigr\}_{h=0}^{H-1}
  \subset s_{t+1:t+H}.
\end{equation}
The recursion is implemented as a plain Python \texttt{for} loop over the
horizon dimension of \texttt{future\_torques}; the prediction from step
$h$ becomes the input to step $h+1$.

\begin{lstlisting}
# MultiStepRobotObjectMCGDFWorldModel.forward (simplified)
state = history_states[:, -1]            # initial s_t
mu, logvar, z, _ = self.encode_context(   # one z per episode (if enabled)
    history_states, history_torques, deterministic_context=...
)
preds = []
for h in range(future_torques.shape[1]):   # H steps
    state, aux = self.dynamics(
        state, future_torques[:, h], damping, friction, object_context, z=z,
    )
    preds.append(state)                    # collect s_{t+h+1}
pred_future = torch.stack(preds, dim=1)    # (B, H, 31)
\end{lstlisting}

\paragraph{No separate ``trajectory decoder''.}
There is deliberately no recurrent network, transformer, or
trajectory-level head that emits $p^o, r^o, v^o, \omega^o$ directly. The
trajectory is generated by $H$ repeated calls to a single neural net
(\texttt{ContactWrenchHead}) embedded inside the analytic Newton--Euler
simulator from Section~\ref{sec:architecture}. All recurrence comes from
plugging $s_{t+h}$ back in as the input to step $h+1$.

\paragraph{The 6-D bottleneck.}
A consequence of Eq.~\eqref{eq:trajectory-recursion} is that the entire
learned signal flowing into the predicted object trajectory at step
$h$ is
\begin{equation}
  F_c(s_{t+h},\,\tau_{t+h+1},\,\xi,\,z) \;\in\; \mathbb{R}^6.
  \label{eq:six-d-bottleneck}
\end{equation}
Everything downstream — Newton's second law for $\dot v^o$, the gyroscopic
Euler equation for $\dot\omega^o$, semi-implicit Euler for
$(p^o, v^o, \omega^o)$, and the exponential-map update for $r^o$ — is
written in PyTorch with no learnable parameters. The object trajectory
therefore depends on the learned model only through these six numbers per
step, repeated $H$ times. This is the structural fix for mismatches
\#1/\#5/\#6/\#7 from the verification report (a single shared
action-reaction wrench, a built-in gyroscopic term, the world-frame
inertia, and no redundant $F_{\mathrm{ext}}$ supervision).

\paragraph{What steers the trajectory.}
The three inputs that effectively determine the predicted object path are
therefore
\begin{itemize}
  \item the \textbf{initial state} $s_t$ (where the cube starts);
  \item the \textbf{future torque sequence} $\tau_{t+1:t+H}$, which drives
    the robot side via Eq.~\eqref{eq:mcgdf-robot-coupled} and, through the
    rolled-out gripper pose, drives $F_c$;
  \item the \textbf{per-episode context} $\xi$ and the latent $z$
    (when the context encoder is enabled), which inform $F_c$ about
    mass, inertia, and material properties.
\end{itemize}

\subsection{Parameter count}

With the default sizes from Table~\ref{tab:subnets}:
\begin{itemize}
  \item DeLaN block (three nets, hidden $=256$, depth $=2$): about
    $219$~k parameters.
  \item Residual head (hidden $=128$, depth $=2$): about $21$~k.
  \item Contact head (hidden $=128$, depth $=3$): about $41$~k.
  \item Joint damping/friction (when learnable): $18$ scalars.
\end{itemize}
Total trainable parameters $\approx 281$~k when
\texttt{learn\_damping\_friction=False} (the default). Most of the budget
is spent on the DeLaN block, which is the heaviest because the
off-diagonal head outputs $36$ values and the inertia heads use the wider
$256$ hidden width inherited from the GT-dynamics baseline.

\section{Forward Kinematics and Jacobian}

The class \texttt{FrankaForwardKinematics} implements the same fixed-link
chain as
\texttt{../GT\_dynamics/eval\_utils.franka\_joint\_and\_gripper\_positions},
fully differentiable in PyTorch.  All seven arm joints are revolute about
the local z-axis, so the geometric Jacobian column for joint $i$ is
\begin{equation}
  J_i \;=\; \begin{bmatrix} z_i \times (p_{\mathrm{ee}} - p_i) \\ z_i \end{bmatrix},
\end{equation}
where $z_i$ is the joint axis and $p_i$ the joint origin, both in the world
frame, and $p_{\mathrm{ee}}$ is the gripper-tip position (hand frame offset
by \texttt{tool\_z\_offset}).  The two finger joint columns are zero,
giving $J_c \in \mathbb{R}^{6\times 9}$ with the last two columns zero.

\section{Contact Wrench Head}

\texttt{ContactWrenchHead} is a small MLP
\begin{equation}
  F_c(s,\tau,\xi) \in \mathbb{R}^6,
  \qquad F_c = (F_c^{\mathrm{lin}},\ F_c^{\mathrm{ang}}).
\end{equation}
It takes the full robot+object state, the recorded torque, and the
per-episode context vector. Its output is interpreted in world frame and
applied to the cube directly; the robot-side reaction torque is derived in
Eq.~\eqref{eq:mcgdf-robot-coupled}.

\section{Residual Head and Pre-Contact Baseline}

The residual head $r_\theta(q,\dot q,\tau)$ is the same as in the robot-only
MCGDF and absorbs the implicit-PD-vs-Python bias plus finite-difference and
numerical residuals.  Because the per-step implicit residual target
\begin{equation*}
  \tau_{\mathrm{ID,GT},t} + d\dot q_t + f\tanh(\dot q_t/\varepsilon_f) - \tau_{\mathrm{applied},t}
\end{equation*}
contains both the implicit-PD bias and the contact wrench effect, the
trainer subtracts the per-episode \emph{pre-contact baseline} of this target
from the per-step target:
\begin{align}
  \mathrm{baseline}_e
  &= \frac{1}{|\mathrm{pre\,contact}|}
     \sum_{t\in\mathrm{pre\,contact}}
     \bigl(\tau_{\mathrm{ID,GT},t} + d\dot q_t + f\tanh(\dot q_t/\varepsilon_f) - \tau_{\mathrm{applied},t}\bigr),
  \\
  r_\theta^{\mathrm{target}}_t
  &= \bigl(\tau_{\mathrm{ID,GT},t} + d\dot q_t + f\tanh(\dot q_t/\varepsilon_f) - \tau_{\mathrm{applied},t}\bigr)
     - \mathrm{baseline}_e.
\end{align}
This is Approach~A of mismatch~\#2.  In effect, the per-episode constant
``non-contact'' part of the residual is removed before $r_\theta$ sees it,
and the contact-induced excess is left for the structured
$\tau_{\mathrm{contact}} = -J_c^\top F_c$ to absorb.

\subsection{Why an internal \texorpdfstring{$F_c$}{F\_c}-based target does
not replace the baseline}
\label{sec:residual-identifiability}

A natural-looking alternative is to skip the pre-contact baseline and
instead define the residual target from the network's own contact-wrench
output:
\begin{equation}
  r_\theta^{\mathrm{target}}_t
  \;=\;
  \underbrace{\bigl(\tau_{\mathrm{ID,GT},t} + d\dot q_t + f\tanh(\dot q_t/\varepsilon_f) - \tau_{\mathrm{applied},t}\bigr)}_{=:\,\rho_t\ \text{(observed residual)}}
  \;+\; J_c(q_t)^\top F_c^{\mathrm{NN}}(s_t,\tau_t,\xi).
  \label{eq:residual-Fc-target}
\end{equation}
The intuition is that when $F_c^{\mathrm{NN}}\!\approx 0$ (pre-contact) the
target reduces to the robot-only residual, and when $F_c^{\mathrm{NN}}$
becomes non-zero (contact) the target self-adjusts. This scheme fails for
the following reason.

\paragraph{Identifiability collapse.}
Substituting Eq.~\eqref{eq:residual-Fc-target} into the supervised loss
gives
\begin{equation}
  \bigl\|r_\theta - r_\theta^{\mathrm{target}}\bigr\|^2
  \;=\;
  \bigl\|(r_\theta - J_c^\top F_c^{\mathrm{NN}}) - \rho_t\bigr\|^2,
  \label{eq:residual-Fc-collapse}
\end{equation}
which is exactly the robot inverse-dynamics constraint: the net learned
residual force on the joints must match the observed residual
$\rho_t$. The existing $\mathcal{L}_{\mathrm{robotGT}}$ and the rollout
MSE already enforce this constraint. Eq.~\eqref{eq:residual-Fc-target}
therefore adds no information about the individual pieces
$(r_\theta,\,F_c)$. The loss is exactly zero everywhere on the hyperplane
\begin{equation}
  \mathcal{H}_t \;=\;
  \bigl\{(r_\theta,\,F_c)\ :\ r_\theta - J_c^\top F_c = \rho_t\bigr\},
\end{equation}
which includes both extreme degeneracies
$(r_\theta=\rho_t,\,F_c=0)$ and
$(r_\theta=0,\,F_c$ does everything$)$, and the symmetry between them is
not broken by any other term in the loss except $\lambda_{\mathrm{contact}}\|F_c\|^2$,
which actively prefers $F_c\!\to\!0$ — i.e.\ the failure mode where
contact is silently absorbed by $r_\theta$.

\paragraph{Bootstrap noise.}
Early in training $F_c^{\mathrm{NN}}$ is random and so is the supervised
target in Eq.~\eqref{eq:residual-Fc-target}. Detaching gradients
(target-network style) keeps the optimization stable but does not
restore identifiability; it only converts fast oscillation into slow
drift along $\mathcal{H}_t$.

\paragraph{Why the baseline works where this does not.}
An identifiability ambiguity between $r_\theta$ and $F_c$ can only be
broken by a label that depends on \emph{one} of the two unknowns. The
pre-contact baseline supplies such a label for $r_\theta$ on a subset of
frames: it uses the external, dataset-level fact that pre-contact frames
have no contact to assert
$r_\theta^{\mathrm{target}}_t \approx 0$ for $t \in \mathrm{pre\,contact}$.
The direct contact-force proxy in
Section~\ref{sec:future-contact-proxy} supplies the symmetric label for
$F_c^{\mathrm{lin}}$. Any scheme that constrains only the combination
$r_\theta - J_c^\top F_c$, including
Eq.~\eqref{eq:residual-Fc-target}, is invariant along $\mathcal{H}_t$ and
cannot separate the two heads.

\section{Loss}

The total loss is
\begin{equation}
  \mathcal{L} =
  \mathcal{L}_{\mathrm{rollout}}
  + \mathcal{L}_{\mathrm{robotGT}}
  + \mathcal{L}_{\mathrm{objectGT}}
  + \mathcal{L}_{\mathrm{DF}}
  + \mathcal{L}_{\mathrm{res,sup}}
  + \mathcal{L}_{\mathrm{reg}}
  + \mathcal{L}_{\mathrm{L2}}.
\end{equation}
\begin{itemize}
  \item $\mathcal{L}_{\mathrm{rollout}}$: weighted MSE on robot state and
    object pose / velocity (sign-invariant quaternion MSE).
  \item $\mathcal{L}_{\mathrm{robotGT}}$: same supervision as the robot-only
    MCGDF for $M, M\ddot q, c, g, \ddot q, \tau_{\mathrm{ID}}$.
  \item $\mathcal{L}_{\mathrm{objectGT}}$: weighted MSE on $\dot v^o$ and
    $\dot\omega^o$ only.  The linearly dependent $F_{\mathrm{ext}}$ target
    is dropped (mismatch~\#7).
  \item $\mathcal{L}_{\mathrm{DF}}$: optional supervision of learnable
    damping/friction scalars against dataset values.
  \item $\mathcal{L}_{\mathrm{res,sup}}$: direct supervision of $r_\theta$
    against the baseline-subtracted target.
  \item $\mathcal{L}_{\mathrm{reg}}$: small L2 regularization on the
    Coriolis term, the residual head magnitude, and the contact wrench
    magnitude.  The contact-wrench regularizer encourages $F_c$ to remain
    near zero on pre-contact windows.
  \item $\mathcal{L}_{\mathrm{L2}}$: explicit L2 penalty on trainable weight
    matrices.
\end{itemize}

\section{Caveats}

\begin{itemize}
  \item \emph{Contact point approximation.}  The Jacobian
    \texttt{FrankaForwardKinematics.jacobian} uses the gripper midpoint as
    the contact location.  Two-finger contacts with distinct contact points
    would require per-finger Jacobians; this is the natural next refinement.
  \item \emph{Friction-cone constraint.}  $F_c$ is unconstrained.  Adding a
    learned friction coefficient and projecting onto the cone
    $\|F_c^{\mathrm{tan}}\|\le\mu\,F_c^{\mathrm{n}}$ would improve
    physical consistency at the cost of training-time non-smoothness.
  \item \emph{Pre-contact baseline assumes a single contact onset per
    episode.}  Re-contacts or multi-phase contact patterns are not currently
    treated separately.  A future variant would window the baseline locally.
  \item \emph{Approach B of mismatch~\#2 (frozen robot-only $r_\theta$).}
    Not yet implemented.  When ready, the trainer would load a robot-only
    MCGDF checkpoint and use its $r_\theta$ output as the supervision target
    here, replacing the per-episode baseline.
  \item The damping double-counting caveat carried over from the robot-only
    MCGDF still applies (see
    \texttt{../../Robot\_only/MCGDF/robot\_only\_mcgdf\_description.tex},
    Caveats); the default \texttt{--omit\_damping} flag avoids it.
  \item \emph{Ground-truth physical context at evaluation
    (train/eval asymmetry).}
    \label{caveat:gt-context-eval}
    The current evaluation scripts
    (\texttt{plot\_multistep\_trajectory.py},
    \texttt{animate\_episode.py}) feed the dataset-recorded
    $\xi = [m_o,\ I_b,\ \mu_o]$ directly into the model at every rollout
    step. $\xi$ enters two places:
    (i)~as part of the \texttt{ContactWrenchHead} input concatenation, and
    (ii)~as the mass and inertia driving the analytic Newton--Euler
    integrator
    (Eqs.~\eqref{eq:mcgdf-object-linear}--\eqref{eq:mcgdf-Iw}).
    Reported gripper / cube prediction errors therefore use \emph{oracle
    physics} and are an upper bound on the accuracy this architecture
    could reach in deployment, not a deployment estimate. The latent
    $z\sim q(z\mid\text{history})$ produced by the
    \texttt{ContextEncoder} is layered on top of $\xi$ rather than
    replacing it. A symmetric ``no oracle at evaluation'' protocol is
    proposed in
    Section~\ref{sec:future-eval-no-gt-context}.
\end{itemize}

\section{Future Works}

\subsection{Direct contact-force proxy supervision}
\label{sec:future-contact-proxy}

The contact wrench head $F_c(s,\tau,\xi) \in \mathbb{R}^6$ is currently
trained without any direct ground-truth label. It is shaped indirectly by
three signals: (i)~the rollout MSE through the
$-J_c^\top F_c$ coupling on the robot side, (ii)~the object linear- and
angular-acceleration supervision through
Eqs.~\eqref{eq:mcgdf-object-linear}--\eqref{eq:mcgdf-object-angular}, and
(iii)~the L2 regularizer $\lambda_{\mathrm{contact}}\|F_c\|^2$ that biases
$F_c$ toward zero in pre-contact windows. PhysX does not expose the
gripper--cube pair impulse through the Isaac~Lab observation API, so the
recorder cannot emit a per-step contact wrench at the time of writing.

The dataset does, however, contain the field
\begin{equation}
  F_{\mathrm{ext,GT},t}
  \;=\;
  \mathtt{object\_dynamics/external\_force\_est\_w}_t
  \;=\;
  m_o\,\bigl(\dot v^o_t - g_w\bigr) \in \mathbb{R}^3,
  \label{eq:future-fext-recipe}
\end{equation}
which is exactly Newton's second law on the cube, solved for the net
external (i.e.\ contact) force, using the finite-difference linear
acceleration $\dot v^o_t$ as recorded. Up to the FD noise inherent in
$\dot v^o_t$, this is a per-step \emph{label} for $F_c^{\mathrm{lin}}$.

\paragraph{Proposed loss term.}
We add a per-step supervised contact-force loss
\begin{equation}
  \mathcal{L}_{\mathrm{contact,sup}}
  \;=\;
  \lambda_{\mathrm{cs}}
  \,\frac{1}{T}\sum_{t=1}^{T}
  \bigl\| F_c^{\mathrm{lin}}(s_t,\tau_t,\xi) - F_{\mathrm{ext,GT},t} \bigr\|_2^2,
  \label{eq:future-contact-sup-loss}
\end{equation}
folded into the total loss alongside the existing
$\mathcal{L}_{\mathrm{objectGT}}$,
$\mathcal{L}_{\mathrm{res,sup}}$, and $\mathcal{L}_{\mathrm{reg}}$ terms.
A new CLI flag \texttt{--lambda\_contact\_sup} (default $0$, opt-in)
controls the weight, mirroring how
\texttt{--lambda\_object\_lin\_acc} and \texttt{--lambda\_contact} are
currently exposed.

\paragraph{Relationship to existing object supervision.}
Eq.~\eqref{eq:future-fext-recipe} and the recorded
$\dot v^o_t$ target used by $\mathcal{L}_{\mathrm{objectGT}}$ are linearly
dependent through the scalar $m_o$:
\begin{equation}
  F_{\mathrm{ext,GT},t} = m_o\,\dot v^o_t - m_o\,g_w.
\end{equation}
Therefore Eq.~\eqref{eq:future-contact-sup-loss} mostly \emph{duplicates}
the lin-acc head of $\mathcal{L}_{\mathrm{objectGT}}$ up to a constant
rescaling and offset. It is not a new independent label. Its practical
value is twofold:
\begin{itemize}
  \item It places the supervision \emph{directly} on
    $F_c^{\mathrm{lin}}$, where gradient flow through the analytic
    Newton--Euler step is removed. This reduces the chain length on
    which the lin-acc loss must rely and, in our experience with the
    robot-only residual head, tends to stabilize early training.
  \item It exposes $\lambda_{\mathrm{cs}}$ as a separate scalar from
    $\lambda_{\mathrm{object\,lin\,acc}}$, letting the practitioner
    weight the two paths independently if FD noise on $\dot v^o_t$ is
    being amplified by the $m_o$ scaling.
\end{itemize}
Because of this redundancy, the recommended default is
$\lambda_{\mathrm{cs}}=0$ with $\lambda_{\mathrm{object\,lin\,acc}}>0$;
the proxy loss is meant as a diagnostic switch, not an additive
improvement.

\paragraph{Caveats.}
\begin{itemize}
  \item \emph{Not a true ground-truth contact force.} Eq.~\eqref{eq:future-fext-recipe}
    is computed from the recorded $\dot v^o_t$, which itself is a
    finite-difference estimate of the cube's world-frame linear
    acceleration. Any FD noise propagates into the label, and that noise
    is amplified by the factor $m_o$ (typically $\mathcal{O}(0.1\,\mathrm{kg})$
    for the cube, but more generally context-dependent through $\xi$).
  \item \emph{No equivalent shortcut for $F_c^{\mathrm{ang}}$.}
    The recorded $\dot\omega^o$ supervises the angular branch only through
    the full gyroscopic equation
    Eq.~\eqref{eq:mcgdf-object-angular}, which depends on $I_w(r^o)$ and
    $\omega^o\!\times\!(I_w\omega^o)$. There is no scalar-rescaled
    label for $F_c^{\mathrm{ang}}$ analogous to
    Eq.~\eqref{eq:future-fext-recipe}. The angular wrench therefore
    remains trained only through the rollout MSE and the angular branch
    of $\mathcal{L}_{\mathrm{objectGT}}$.
  \item \emph{Pre-contact regime.} On pre-contact frames, both
    $F_{\mathrm{ext,GT},t}$ and the target for $F_c^{\mathrm{lin}}$ are
    approximately zero (up to gravity already removed in the recipe).
    Adding Eq.~\eqref{eq:future-contact-sup-loss} on pre-contact windows
    is therefore consistent with $\lambda_{\mathrm{contact}}\|F_c\|^2$
    and does not require a separate masking schedule.
\end{itemize}

\subsection{Alternative routes for a true contact-wrench label}

The proxy in Eq.~\eqref{eq:future-fext-recipe} is bounded above by the FD
quality of $\dot v^o_t$ and cannot supervise $F_c^{\mathrm{ang}}$. Two
heavier-touch options would provide a genuine, low-noise label:
\begin{itemize}
  \item \emph{Extend the recorder to log PhysX
    \texttt{body\_incoming\_joint\_wrench\_b}.} This API returns the
    wrench transmitted through an articulation joint. For the cube it is
    of limited use because the cube is a \emph{free body}, not an
    articulation link, and so the wrench would be reported as zero. The
    useful application of this signal is on the Franka side: querying
    the wrench at the hand--finger joints exposes the contact wrench in
    a body frame, which after a known finger-to-world rotation gives a
    direct label for the wrench the gripper applies to the cube
    (Newton's third law gives $F_c$ in world frame).
  \item \emph{Add a custom Isaac~Lab observation that reads the PhysX
    contact buffer directly.} PhysX maintains per-pair contact impulses
    at the contact-point level; aggregating impulses over all
    cube--gripper contact pairs and dividing by $\Delta t$ yields a
    per-step ground-truth $F_c$ (and, with the contact-point offset,
    a torque about the cube center). This is the most invasive option,
    requiring a new observation term in
    \texttt{robot\_object\_mcgdf\_env\_cfg.py} and a new field in
    \texttt{RobotObjectMCGDFRolloutDataset}, but it produces the only
    truly unambiguous label for the full $F_c \in \mathbb{R}^6$.
\end{itemize}
Either route would replace Eq.~\eqref{eq:future-contact-sup-loss} with a
direct $\|F_c - F_c^{\mathrm{GT}}\|^2$ term and remove the FD-noise
caveat.

\subsection{Coupling-driven object dynamics with an online latent context update}
\label{sec:future-eval-no-gt-context}

This subsection replaces an earlier sketch in which the explicit
physical context $\xi = [m_o,\ I_b,\ \mu_o]$ was estimated by a
regression head $\hat\xi(z)$ and substituted at evaluation. Two
practitioner-side concerns rule that earlier sketch out:
\begin{enumerate}
  \item Forcing $z$ to predict the dataset's $\xi$ pins a fixed semantic
    role on the latent (``the model's guess at $\xi$'') and adds
    parameters (a Cholesky-parameterized $3\times 3$ inertia and a
    softplus mass head). The latent should be free to encode whatever
    residual context is useful, not a copy of the recorded labels.
  \item Treating the latent as a single one-shot encoding of an initial
    history window throws away later contact evidence. The
    instantaneous discrepancy between the model's prediction and the
    actually observed cube motion is exactly the signal that ought to
    refine the latent over time.
\end{enumerate}
Two structural changes are therefore proposed.

\paragraph{Change 1 -- object state predicted from the robot's ODE
plus a contact-gated coupling, instead of from Newton--Euler.}
The free-rigid-body integrator
Eqs.~\eqref{eq:mcgdf-object-linear}--\eqref{eq:mcgdf-quat-exp} is removed
from the object side. The object substate
$(p^o, r^o, v^o, \omega^o)$ is instead driven by the gripper kinematics
through a learned, state-dependent coupling factor
$\alpha_t \in [0,1]$ that gates ``how attached'' the cube is to the
gripper at the current step. Concretely, let
\begin{equation}
  \mathbf{d}_t \;=\; p_g(q_t) - p^o_t \;\in\;\mathbb{R}^3,
  \qquad
  \alpha_t \;=\; \sigma\!\bigl(\mathrm{MLP}_\alpha(\|F_c^{\mathrm{lin}}\|,\,\mathbf{d}_t,\,z_t)\bigr),
  \label{eq:future-coupling-alpha}
\end{equation}
where $p_g(q)$ is the analytic gripper-tip position from
\texttt{FrankaForwardKinematics} and the gripper twist
$(v_g, \omega_g) = J(q)\dot q$ is the existing geometric Jacobian
product. The cube state is then advanced by blending the gripper twist
into the cube twist:
\begin{align}
  v^o_{t+1}      &= \alpha_t\,v_g(q_t,\dot q_t)
                   \;+\; (1-\alpha_t)\,\rho_v\,v^o_t,
                   \label{eq:future-vobj}\\
  \omega^o_{t+1} &= \alpha_t\,\omega_g(q_t,\dot q_t)
                   \;+\; (1-\alpha_t)\,\rho_\omega\,\omega^o_t,
                   \label{eq:future-womega}\\
  p^o_{t+1}      &= p^o_t \;+\; v^o_{t+1}\,\Delta t,
                   \label{eq:future-pobj}\\
  r^o_{t+1}      &= r^o_t \otimes \exp\!\bigl(\tfrac{1}{2}\omega^o_{t+1}\,\Delta t\bigr),
                   \label{eq:future-robj}
\end{align}
where $\rho_v,\rho_\omega\in(0,1]$ are small learned decay scalars
(softplus-bounded) that model passive ``rest on the table'' losses.
Intuition: when the gripper is close to the cube and pressing on it,
$\alpha_t\!\to\!1$ and the cube inherits the gripper's twist; when out
of contact, $\alpha_t\!\to\!0$ and the cube's twist relaxes to zero so
$p^o$, $r^o$ stay put. No $m_o$, $I_b$ or $\mu_o$ appears anywhere in
Eqs.~\eqref{eq:future-vobj}--\eqref{eq:future-robj}.

\paragraph{What is gained, and what is lost, by Change 1.}
\emph{Gained.} The object-side parameter budget collapses to one small
MLP for $\alpha_t$ and two scalars $\rho_v,\rho_\omega$. The dataset's
$\xi$ never enters the model's forward pass at any time, so there is
no GT-at-eval question to answer: $\xi$ becomes purely a label used by
training-time losses. The inductive bias ``the cube only moves when
carried'' matches the Franka Lift task structurally and is robust
out-of-distribution as long as that dichotomy holds.\\
\emph{Lost.} The object side no longer obeys Newton--Euler. Cases the
new integrator cannot represent include ballistic motion (a thrown
cube), free spinning after release, motion driven by an external
impulse other than the gripper, and any task where the cube's mass or
inertia visibly governs its trajectory. The 6-D wrench $F_c$ now
enters Eq.~\eqref{eq:future-coupling-alpha} only through
$\|F_c^{\mathrm{lin}}\|$ as a gating signal; it no longer constrains
the cube's acceleration via $F = m a$. For the Lift task this trade is
acceptable, but it should be flagged for any task that involves true
free-flight phases.

\paragraph{Change 2 -- online latent update from the per-step
prediction residual.}
The variational $z$ is currently produced once at the start of the
rollout from a $K$-step history window and held fixed for the full
horizon $H$. The proposed change is to make $z$ \emph{evolve over
time}, refined at every step by the one-step prediction residual:
\begin{equation}
  e_t \;=\; \mathrm{Proj}_{\mathrm{obj}}\!\bigl(s_t^{\mathrm{obs}} - \hat s_t\bigr)
  \;\in\;\mathbb{R}^{13},
  \qquad
  z_t \;=\; z_{t-1} \;+\; \mathrm{GatedUpdate}\!\bigl(z_{t-1},\,e_t,\,F_c^{t-1}\bigr).
  \label{eq:future-latent-online-update}
\end{equation}
Here $\hat s_t$ is the model's one-step prediction made under
$z_{t-1}$, $\mathrm{Proj}_{\mathrm{obj}}$ selects the object substate
($p^o, r^o, v^o, \omega^o$) from the residual since that is the part
the contact dynamics can do something about, and
$\mathrm{GatedUpdate}$ is a small GRU-style cell of width
$\mathrm{latent\_dim}$ -- think of it as the analog of a Kalman
innovation update on $z$. The initial value $z_0$ is still produced by
the existing \texttt{ContextEncoder} from the warm-up window so the
filter has a sensible prior to start from.

Conceptually this turns $z$ from a one-shot encoding into a sequential
posterior estimate: each frame of newly-observed contact behaviour
nudges $z$ in the direction the model needed in order to have been
correct, and the next $F_c$ and $\alpha$ benefit from that correction
without any additional supervised label.

\paragraph{Training-time protocol.}
At training time, the closed-loop residual signal is always available
(the dataset supplies $s_t^{\mathrm{obs}}$). The trainer therefore
runs the model in \emph{closed-loop on a sliding window} of length
$L_{\mathrm{filter}}$:
\begin{lstlisting}
z = ContextEncoder(history[:K])              # prior from warm-up window
for t in range(K, K + L_filter):             # closed-loop filter window
    s_hat = OneStep(s_obs[t-1], tau[t], z)
    e     = Proj_obj(s_obs[t] - s_hat)
    z     = z + GatedUpdate(z, e, F_c)
for h in range(rollout_horizon):             # open-loop rollout window
    s = OneStep(s, tau[K + L_filter + h], z) # z is frozen here
\end{lstlisting}
The supervised losses (rollout MSE, robot- and object-dynamics
supervision, residual head with pre-contact baseline, KL on $z_0$)
operate on the rollout-window predictions versus
$s^{\mathrm{obs}}_{K+L_{\mathrm{filter}}:K+L_{\mathrm{filter}}+H}$ as
today. Truncated BPTT through the filter recurrence over
$L_{\mathrm{filter}}$ steps keeps the backward pass affordable.
The GT context $\xi$ is now used \emph{only} by the supervised losses
(InfoNCE anchor on $z_0$, object-dynamics targets, residual-baseline
computation); it never enters the model's forward pass.

\paragraph{Evaluation-time protocol.}
The same protocol applies, with one critical observation: during the
filter window the model is permitted to read the dataset's
$s_t^{\mathrm{obs}}$ because that information is also available in
deployment from cheap, GT-context-free sensors (joint encoders plus
the recorded cube pose from an onboard RGB-D pipeline). It is not
``oracle context.'' The open-loop rollout window then uses the
filter-refined $z_{K+L_{\mathrm{filter}}}$ and predicts forward without
peeking. Evaluation scripts gain a single flag
\texttt{--filter\_steps} controlling $L_{\mathrm{filter}}$; the
existing \texttt{start\_t} and \texttt{rollout\_steps} continue to
control the rollout window.

\paragraph{Open design choices to discuss.}
\begin{itemize}
  \item \emph{Object-only vs.\ full residual in $e_t$.} Restricting $e_t$
    to the object substate (as written) focuses the filter on the
    contact-coupled coordinates and avoids being driven by
    joint-tracking noise on the robot side. Using the full $31$-vector
    is simpler but noisier; this is a small ablation.
  \item \emph{Single $\alpha_t$ vs.\ two-component $(\alpha^{\mathrm{rigid}}, \alpha^{\mathrm{slide}})$.}
    A single scalar conflates two regimes the dataset actually contains:
    ``cube rigidly held by closed gripper'' and ``cube sliding under
    pushing contact.''  A two-headed $\alpha_t$ -- one rigid, one tangential --
    would generalise to in-plane pushing without sacrificing the lift
    inductive bias.
  \item \emph{Keep $F_c \in \mathbb{R}^6$ or replace by a scalar
    contact head?} $F_c$ no longer drives the object integrator, only
    the gating. We recommend keeping the 6-D shape so the robot-side
    coupling $-J_c^\top F_c$ (Newton's third law on the robot, the
    structural reason the robot+object split exists) is preserved and
    the residual head $r_\theta$ does not have to absorb the contact
    reaction torque.
  \item \emph{Where the warm-up $K$ and filter $L_{\mathrm{filter}}$
    are reported.} The reported open-loop horizon-$H$ error becomes a
    function of $(K, L_{\mathrm{filter}}, H)$; a sweep over
    $L_{\mathrm{filter}}$ is the natural plot to show the value of the
    online filter -- error should drop monotonically with
    $L_{\mathrm{filter}}$ until the latent saturates.
  \item \emph{Relation to the project proposal's Pi-CDD year-1.} The
    filter recurrence Eq.~\eqref{eq:future-latent-online-update} is
    exactly the time-evolving ``Context $c_t$'' channel sketched in the
    proposal's Figure 1, with the difference that $c_t$ here is updated
    by an innovation rather than by re-encoding a sliding window. The
    coupling formulation Eqs.~\eqref{eq:future-vobj}--\eqref{eq:future-robj}
    fits the proposal's ``Action--Reaction'' framing because $\alpha_t$
    is precisely the model's read of how strongly the gripper action is
    coupled to the object reaction at the current step.
\end{itemize}

\paragraph{What this looks like at the code level (sketch).}
\begin{itemize}
  \item \emph{models.py.} Replace \texttt{RobotObjectMCGDFStep.\_object\_dynamics}
    with the coupling block in
    Eqs.~\eqref{eq:future-coupling-alpha}--\eqref{eq:future-robj}.
    Add a small \texttt{CouplingHead} (MLP) and a
    \texttt{LatentInnovationUpdate} (GRUCell-style) module. The DeLaN
    block, residual head, contact-wrench head, FK chain, and quaternion
    exp map are unchanged.
  \item \emph{wm\_train.py.} Add \texttt{--filter\_steps} CLI flag.
    Modify \texttt{run\_epoch} to run the closed-loop filter window
    before the open-loop rollout window; the supervised losses keep
    their current API.  Optionally drop or down-weight the existing
    KL-only-on-$z_0$ term in favour of a small KL prior on the
    sequence $(z_K, z_{K+1}, \dots)$.
  \item \emph{plot\_multistep\_trajectory.py / animate\_episode.py.}
    Add \texttt{--filter\_steps} and route the filter window through
    the existing \texttt{model.encode\_context} +
    \texttt{LatentInnovationUpdate} path. No GT $\xi$ is loaded for
    the model forward pass; the regression-head substitution sketched
    in the previous (superseded) plan is not needed.
\end{itemize}

\paragraph{Status of the previous plan.}
The earlier sketch of training a 13-dim $\hat\xi(z)$ regression head
and substituting $\hat\xi$ at evaluation is dropped; this subsection
replaces it.

\subsection{Other planned extensions}

\begin{itemize}
  \item \emph{Approach~B of mismatch~\#2 (frozen robot-only $r_\theta$
    as supervision).} Train a robot-only MCGDF first, freeze its
    $r_\theta$ output, and use it as the supervision target for the
    residual head here, replacing the per-episode pre-contact baseline.
    This removes the assumption of a single contact onset per episode
    and would be the natural pairing with the per-pair contact-buffer
    route above.
  \item \emph{FiLM-modulated DeLaN heads.} Inject the per-episode
    context $\xi$ into the DeLaN block via FiLM layers, matching the
    treatment of $r_\theta$ and $F_c$. The current implementation
    leaves $H,c,g$ context-free because the robot is fixed across
    episodes; this remains correct when only the cube changes, but
    becomes necessary once payload or in-hand objects affect the
    effective robot inertia.
  \item \emph{Per-finger Jacobians and friction cone.} Replace the
    single-midpoint Jacobian with per-finger Jacobians and constrain
    $F_c$ to lie inside a learned friction cone
    $\|F_c^{\mathrm{tan}}\|\le \mu(\xi)\,F_c^{\mathrm{n}}$. This pairs
    naturally with the explicit contact-buffer label, which already
    provides per-contact normals.
\end{itemize}

\end{document}
